
\chapter{État de l'art en modélisation de bâtiments : courte étude bibliographique}
    \section{Analogie avec les circuits électriques}~\cite[chap.~VI]{manuel_energie_batiment}
    \section{Modèle de Contrôle Prédictif pour le bâtiment}~\cite[chap.~I--IV]{MPC_needs}
%        \begin{figure}
%            \centering
%            \begin{tikzpicture}[node distance=1.5cm]
%
%                % Nodes
%                \node (timek) [startstop] {@ temps $t_k$};
%                \node (measurements) [process, below of=timek] {Prendre les mesures du processus};
%                \node (model) [process, below of=measurements, minimum width=8cm] {%
%                    \textbf{Modèle du processus} \\[1ex]
%                    \begin{minipage}{7cm}
%                        \textbf{Entrées :} \\
%                        - Actions de contrôle actuelles et futures \\
%                        - Perturbations \\[1ex]
%                        $\Rightarrow$ \textbf{Sorties futures du processus}
%                    \end{minipage}
%                    \hspace{0.5cm}
%                    \begin{minipage}{3cm}
%                        \textbf{Objectifs} \\
%                        \textbf{Contraintes}
%                    \end{minipage}
%                };
%                \node (optimization) [process, below of=model] {Résoudre le problème d'optimisation ci-dessus \\ $\Downarrow$ Meilleures actions de contrôle actuelles et futures};
%                \node (implement) [process, below of=optimization] {Mettre en œuvre la meilleure action de contrôle \textit{actuelle}};
%                \node (timek1) [startstop, below of=implement] {temps $t_{k+1}$};
%
%                % Arrows
%                \draw [arrow] (timek) -- (measurements);
%                \draw [arrow] (measurements) -- (model);
%                \draw [arrow] (model) -- (optimization);
%                \draw [arrow] (optimization) -- (implement);
%                \draw [arrow] (implement) -- (timek1);
%                \draw [arrow] (timek1.east) .. controls +(2,0) and +(2,0) .. (timek.east);
%
%            \end{tikzpicture}
%            \caption{Schéma résumant le principe du MPC (\cite{MPC_needs})}
%            \label{fig:mpc}
%        \end{figure}

    \section{Machine Learning pour le bâtiment}
        \subsection{Approches traditionnelles}
            \cite{WML}
        \subsection{Réseaux de neurones}
            \cite{LeCun_DL}
            \cite{UDL}